\chapter[Goetsu font test]{Goetsu font test}

% ---------------------------------------------------------------------------
% This chapter exercises Jyutcitzi vs Goetsu routing.
% Body text runs under \jczOn (see narihato-orfu-analects.tex).
% Goetsu blocks suspend Jyutcitzi, typeset in goetsusioji, then restore.
% ---------------------------------------------------------------------------

\section*{Raw PUA (global \texttt{\textbackslash jczOn})}



\section*{Explicit \texttt{\textbackslash jcz\{...\}}}

\jcz{}

\section*{Goetsu environment}

\begin{goetsu}

\end{goetsu}

\section*{Goetsu inline}

漢字\goetsuinline{}漢字

\section*{Goetsu in chapter title (real use case)}

See the running head / TOC entry for this chapter, and the title line above:
Goetsu glyphs can sit inside \verb|\chapter{...\begin{goetsu}...\end{goetsu}}|.

\section*{Mixed body + Goetsu block}

請問吳語的副詞、副動詞的系統閣下有過闡述嗎？比如 uli、suli、syli、deuli 等等這些
\begin{goetsu}

\end{goetsu}

\qaanswer{%
Placeholder answer for layout test only.
}
